% 摘要章节
% 说明：
% 1. 本文件设计为被主文档 \input{} 进来使用。
% 2. 内容基于论文的引言、方法与实验结果撰写。

\begin{abstract}

增量式三维场景重建要求系统在逐帧接收新观测的同时持续扩展地图，而有限的在线优化预算使得新插入高斯的初始属性质量直接影响地图的累积效果。现有基于 3D Gaussian Splatting 的增量式建图系统在表示与初始化两方面仍存在不足：三维椭球原语缺乏显式的表面法线语义，基于规则的初始化策略也难以充分利用外部几何先验。针对这两个相互关联的问题，本文提出了一种面向增量式 2DGS 场景重建的上下文感知前馈高斯回归系统。系统以 2DGS surfel 为地图基础表示，融合 SPNet 深度补全与单目法线先验构造候选点及其几何方向；前馈回归网络通过跨注意力机制聚合当前帧与历史帧的图像--渲染上下文，并以法线先验为锚点经 Gram-Schmidt 正交化构造 surfel 旋转，在单次前向传播中直接预测完整的 surfel 属性。网络训练通过自蒸馏闭环完成，无需外部标注数据集。在 R3LIVE 和 MCD 数据集上的实验表明，本文方法在全部测试序列上均优于现有基线，较最接近的多传感器增量高斯建图基线平均提升 1.28\,dB PSNR；消融实验进一步验证了前馈回归器在有限优化预算下对新高斯早期收敛质量的稳定改善。

\end{abstract}
